\title{LongRoPE: Extending LLM Context Window Beyond 2 Million Tokens}

\begin{abstract}
Expanding the context window in large language models (LLMs) is highly advantageous. Nonetheless, present approaches to increasing the window size face several obstacles, such as the significant expense of fine-tuning, the limited availability of lengthy text data, and disruptive effects caused by introducing unfamiliar token positions. As a consequence, the maximum effective context window currently remains near 128k tokens.
\end{abstract}

This work presents LongRoPE, a novel approach that, for the first time, enables pre-trained LLMs to process context windows as large as \textbf{2048k} tokens, requiring no more than 1k fine-tuning steps at training lengths up to 256k, all while preserving performance on the original, shorter contexts. LongRoPE accomplishes this through three main contributions: (i) we detect and leverage two distinct non-uniform patterns in positional interpolation via an efficient search procedure, resulting in superior initialization for fine-tuning and permitting an 8$\times$ context length increase without additional fine-tuning; (ii) we develop a two-phase progressive extension method, initially fine-tuning on 256k tokens before applying a second positional interpolation to the already extended model to reach the 2048k context window; (iii) we realign LongRoPE at the 8k context length to recover performance for shorter windows. Comprehensive evaluations using LLaMA2 and Mistral on diverse benchmarks validate our approach. Models adapted with LongRoPE preserve the underlying architecture with only minor positional embedding alterations, ensuring compatibility with existing optimizations. Code will be released at \texttt{https://github.com/microsoft/LongRoPE}
\end{abstract}

\section{Introduction}

Large Language Models (LLMs) have achieved impressive results across a range of benchmarks~\cite{gpt4,llama2}, yet they are constrained by relatively short context windows, such as the 4096-token boundary in LLaMA2~\cite{llama2}. When prompted with inputs exceeding this limit, LLMs exhibit degraded performance, as the models are not trained to handle content in those additional positional encodings. This constraint creates difficulties in scenarios that demand processing extensive information within context, for instance, in-context learning requiring many examples~\cite{Huang2023FewerIM} or in deploying LLM agents~\cite{park2023generative,madaan2023selfrefine}.

Recent studies demonstrate that the context window of a pre-trained LLM can be enlarged up to approximately 128k tokens through fine-tuning on longer sequences~\cite{longlora,pi,yarn,zhang2024soaring,liu2023scaling}. Extending the context window further, however, is impeded by three primary challenges. Firstly, introducing new position indices without training causes significant distributional anomalies and severe convergence issues during fine-tuning, due to the prevalence of anomalous values~\cite{pi}. This difficulty becomes especially pronounced in cases where the context window is expanded from 4k to over $>$1000k tokens, resulting in the addition of more than 90\% new, untrained positions. Secondly, successful fine-tuning presupposes access to texts with adequate lengths. Unfortunately, currently available datasets rarely provide examples exceeding 1000k tokens, and employing such lengthy data is computationally demanding, often requiring excessive GPU usage and lengthy training times. Finally, pushing the context window to extreme lengths causes the attention mechanism to be diluted among the multitude of token positions, which undermines the model's effectiveness on tasks involving original, short context lengths~\cite{pi}.

A common strategy to address the initial issue involves interpolating RoPE positional embeddings~\cite{rope,pi}, which essentially adjusts new positional indices to fit within the bounds established during pretraining, as illustrated in Fig.\ref{fig:overview}. The Position Interpolation (PI) technique~\cite{pi} performs linear interpolation of RoPE's rotary angles based on the expansion ratio. The NTK method~\cite{ntk,dynamicntk} proposes non-uniform schemes for interpolation and extrapolation along the dimensions of RoPE. Meanwhile, YaRN~\cite{yarn} separates RoPE dimensions into three distinct frequency-based clusters, employing extrapolation, NTK, and linear interpolation for each group accordingly. Nonetheless, positional embedding within the Transformer reveals a \emph{complex, uneven distribution of information entropy}. Current methods do not sufficiently exploit these intricate, irregular patterns, resulting in a loss of valuable information and ultimately constraining the maximum achievable context window.

Section~\ref{sec:analysis} provides empirical evidence for two principal conclusions: \textbf{(1)} Successful positional interpolation must address two distinct kinds of non-uniformities—those related to differences across RoPE dimensions and those stemming from varying token positions. Specifically, smaller RoPE dimensions and the earliest token positions are best served by reduced interpolation. However, the optimal interpolation configuration is also a function of the intended maximum sequence length for extension. \textbf{(2)} Integrating these observed non-uniformities into the interpolation strategy enables the model to more faithfully preserve crucial information present in the original RoPE embeddings, notably within key dimensions and for important token positions. This approach curtails informational loss due to interpolation, thereby yielding superior initialization for subsequent fine-tuning stages. Furthermore, this method supports sequence length extension by up to 8$\times$ even in the absence of fine-tuning.

Building upon these results, we present \textbf{LongRoPE}, a powerful approach that increases the LLM context window to over 2 \emph{million} tokens. The foundation of LongRoPE lies in three principal advancements. Firstly, LongRoPE leverages multidimensional variations in positional interpolation to their fullest extent. By assessing token locations, LongRoPE determines optimal rescale factors for the rotation angles used by RoPE across all its dimensions. Considering that the space required to explore suitable rescale factors grows exponentially as the extension ratio increases, LongRoPE employs an evolutionary search method enhanced with two additional optimization strategies to maximize efficiency in this process. As illustrated in Fig.~\ref{fig:overview}, the result of the search is an adapted RoPE configuration with rescaled parameters.

Subsequently, {LongRoPE} adopts a resourceful, stepwise extension approach to reach a context window of 2048k, circumventing the requirement for direct fine-tuning with ultra-long text sequences, which are uncommon and challenging to obtain. The process initiates by identifying a suitable 256k length for the pre-trained LLM, followed by fine-tuning at that length. Leveraging the capability of non-uniform positional interpolation to permit up to an 8$\times$ expansion in settings where fine-tuning is not necessary, we proceed with an additional search for novel RoPE rescale factors on the previously fine-tuned, lengthened LLM. This process effectively enables a 2048k context window for models such as LLaMA2 and Mistral~\cite{mistral}.

In order to prevent loss of performance within the original, more limited context window, {LongRoPE} maintains dynamic adjustment of the RoPE rescale parameters on the expanded LLM. As with the process of increasing the window from 256k to 2048k, we also decrease the scale to 4k and 8k context windows for the LLM fine-tuned at 256k, leveraging our search method to favor reduced positional interpolation. For inference cases where the input sequence is below 8k tokens, RoPE is updated utilizing the optimized rescale factors identified by the search procedure.

Through comprehensive experimentation involving multiple LLMs and a range of long-context tasks, our approach demonstrates significant efficacy. Our results indicate that LongRoPE consistently preserves low perplexity across evaluation lengths from 4k up to 2048k, surpasses 90\% accuracy in passkey retrieval, and maintains performance comparable to existing baselines on standard benchmarks situated within a 4096-token context window. Moreover, LongRoPE is broadly applicable to all LLMs utilizing RoPE embedding. The codebase and the LongRoPE-2048k models will be made publicly available.

\section{Non-uniformity in Positional Interpolation}
\label{sec:analysis}

\subsection{Preliminary}

Transformer models require explicit positional information, often in the form of position embedding, to represent the order of input tokens. Our work focuses on the RoPE~\cite{rope} position embedding, which is widely used in recent  LLMs. For a token at position index $n$, its corresponding RoPE encoding can be simplified as follows:

\begin{equation}
		\fontsize{7.8}{7.8} \selectfont
	\label{eq:rope}
	\left[ cos(n\theta_0), sin(n\theta_0), cos(n\theta_1), \cdot\cdot\cdot, cos(n\theta_{d/2-1}), sin(n\theta_{d/2-1}) \right]   
\end{equation}

Here, $d$ denotes the dimensionality of the embeddings,
$n\theta_i$ specifies the rotation angle assigned to the token located at the $n$-th position,
and $\theta_i=\theta^{-2i/d}$ characterizes the set of rotation frequencies. Within RoPE, the standard base for $\theta$ is set to 10000.

\textbf{\emph{Context window extension ratio $s$} and positional interpolation}. The parameter $s$ is introduced as the proportion between the new, extended context window length $L'$ and the original context window length $L$, formally given by $s=\frac{L'}{L}$.

To extend the context window from $L$ to $L'$, current positional interpolation methods suggest downscaling rotation frequencies $\theta_i$ based on extension ratio $s$. Let $\beta=\theta^{2/d}$, and $\lambda$ denote the actual rescale factor related to $s$, we   unify these positional interpolation methods as follows:

\begin{equation}	
	\fontsize{7.5}{7.5} \selectfont
	\label{eq:unified_rope}
	\begin{aligned}
		\left[cos\left(\frac{n}{\lambda(\beta)^0}\right), sin \left(\frac{n}{\lambda(\beta)^0}\right), cos\left(\frac{n}{\lambda(\beta)^1}\right), \cdot\cdot\cdot,  sin \left(\frac{n}{\lambda(\beta)^{d/2-1}}\right) \right]   
	\end{aligned}
\end{equation}

\textbf{Linear positional interpolation (PI)}. PI~\cite{pi} proposes interpolating the position indices linearly for sequences that fall within the pre-trained context window. When extending to a new sequence length with ratio $s$, the method scales down the rotation angles for each position uniformly by $\lambda=s$ over every RoPE dimension. This approach, however, causes the representation of positional information to become compressed, making it difficult for the model to differentiate tokens with neighboring positions. Consequently, the effectiveness of PI declines when the extension ratio is large.

\textbf{NTK-based interpolation and extrapolation}. Studies such as~\cite{ntk,dynamicntk} examine RoPE through the lens of information encoding and utilize Neural Tangent Kernel (NTK) theory~\cite{jacot2018neural,tancik2020fourier} to inform their approach. To address the issue of position overcrowding inherent in PI, these works propose distributing the effects of interpolation more evenly across the different RoPE dimensions. By moderating the scaling so that dimensions corresponding to lower frequencies are scaled up more than those representing higher frequencies, they achieve both interpolation and extrapolation capabilities, parameterized as $\lambda=s^{i}$. The dynamic NTK method introduced in~\cite{dynamicntk} further refines this process by adapting the extension ratio at each location in response to the observed sequence length. Distinct from PI, which depends on additional fine-tuning, NTK-informed strategies facilitate extending the context window without fine-tuning, typically permitting up to a 4$\times$ extension.

\textbf{YaRN}~\cite{yarn} proposes a notable enhancement for positional interpolation by categorizing RoPE dimensions according to their associated frequencies and applying distinct interpolation methods to each group. Specifically, it assigns extrapolation ($\lambda$=1) to the highest frequency dimensions, implements linear interpolation (PI) for the lowest frequencies, and applies the NTK interpolation for those dimensions that lie between. The main innovation of YaRN is its frequency-based partitioning of RoPE dimensions, a process that currently relies on manual, empirical determination. Consequently, this approach may not yield optimal outcomes when adapting to novel LLMs.

\subsection{Study on Non-uniform Positional Interpolation}

Drawing on insights from NTK and YaRN, we observe their improvements are largely due to the incorporation of non-linearity, particularly by analyzing diverse frequencies along the RoPE dimensions for tailored interpolation and extrapolation. Yet, existing non-linear approaches are predominantly based on manually crafted heuristics. This leads to two critical questions: (1) Can positional interpolation be further optimized beyond its present form? (2) Are there novel non-linearities that have not yet been investigated?

In order to address these inquiries, we employ evolution search (refer to Sec.~\ref{sec:method}) to identify improved non-uniform positional interpolations tailored to LLaMA2-7B. The process utilizes perplexity as the optimization criterion and selects 5 random instances from the PG19~\cite{pg19} validation dataset. Our experimental evaluation yields several significant insights.

\textbf{Finding 1}: \textit{RoPE dimensions exhibit substantial non-uniformities, which are not effectively handled by current positional interpolation methods.}

We determine the optimal value of $\lambda$ for each RoPE dimension using Eq.~\ref{eq:unified_rope}. In Table~\ref{tbl:exp1}, we present a comparison of perplexity scores for LLaMA2-7B across various approaches, evaluated on the PG19 and Proof-pile~\cite{proof-pile} test sets, all without any additional fine-tuning. Our search-derived configuration yields marked enhancements over existing methods, indicating that both the linear (PI) and non-uniform (Dynamic-NTK and YaRN) interpolation strategies are not ideal. Of particular interest, YaRN achieves lower performance than PI and NTK on PG19, attributed to its failure to accommodate the target context window length in LLMs without fine-tuning; for instance, YaRN’s perplexity escalates rapidly beyond 7k tokens when the context size is set to 8k.

Our investigation reveals that the scaling coefficients $\lambda$ in Eq.~\ref{eq:unified_rope} exhibit variability, contrasting with the constant scale $s$ employed in PI, the formulaic approach of NTK, and the group-based computation in YaRN. These varied scaling factors yield notable gains in LLaMA2's language modeling accuracy (measured by perplexity) for extended context lengths of 8k and 16k tokens, even in the absence of additional fine-tuning. The underlying reason is that the revised positional embeddings retain the essential structure of RoPE—particularly in key dimensions—which helps the LLM more accurately differentiate between tokens in proximal positions.

\textbf{Finding 2}: \textit{RoPE for the initial tokens in the input sequence should be extrapolated with less interpolation.} 

We propose that the RoPE interpolation for the first $\hat{n}$ tokens in input sequences should be minimized, as these tokens are typically assigned higher attention scores and play a pivotal role in attention mechanisms, an effect documented by Streaming LLM~\cite{streamingllm} and LM-Infinite~\cite{han2023lminfinite}. To test this premise, we expand the context window to 8k and 16k using PI and NTK methods, selectively disabling interpolation for the initial $\hat n$ tokens, where $\hat n$ varies across (0, 2, ..., 256). Setting $\hat n$ to zero recovers the default PI and NTK configurations. Table~\ref{tbl:tokenposition} provides evidence of two principal findings: \textbf{(1)} the omission of position interpolation for early tokens results in measurable performance gains; \textbf{(2)} the best value of $\hat n$ is influenced by the target context window size.

\textbf{Finding 3}: \textit{Non-uniform positional interpolation effectively extends LLM context window in both fine-tuning and non-fine-tuning settings.} 

Although our explored non-uniform position interpolation already yields notable gains in extension performance at 8k and 16k context lengths without additional training, extending to longer sequences necessitates fine-tuning. Therefore, we fine-tune LLaMA2-7B with our optimized RoPE approach for a context window of 64k tokens (fine-tuning settings detailed in the Appendix). As illustrated in Table~\ref{tbl:finetune}, this technique outperforms both PI and YaRN methods by a wide margin, whether evaluated pre- or post-fine-tuning. The improvement arises from the deliberate application of non-uniform positional interpolation, which reduces information loss and establishes superior initial parameters for fine-tuning.

\textbf{Summary}. In this work, we identify and leverage two forms of non-uniformity—across RoPE dimensions and token positions. Employing these non-uniform characteristics in positional interpolation substantially enhances the performance of large language models when expanding their context lengths.

\section{LongRoPE}

\label{sec:method}
Building on these insights, we propose LongRoPE, which initially leverages a streamlined search algorithm designed to harness both identified non-uniformities. Utilizing this approach, we subsequently expand the LLM context window to over 2 million tokens.

\subsection{Problem Formulation}

These two sources of non-uniformity expand the range of possible solutions and complicate the optimization process. To manage these challenges, we recast the multidimensional non-uniform position interpolation optimization task as a search problem.

For a LLM targeting a  context window size of $L'$ and lengthy input documents $\mathbf{X}$, where each $\mathbf{x}\in\mathbf{X}$ surpasses $L'$ in token length, we denote the original rotary angle of the $i^{th}$ dimension in RoPE embedding at token position $n$ as  $\frac{n}{\beta^i}$. The optimization problem is then formulated as follows:

\begin{equation}
\begin{gathered}
		\label{eq:problem}

\underset {\mathbf{x}\in\mathbf{X}; \, |\mathbf{x}|\ge L'}{ \operatorname {arg\,min} } \,  \mathcal{L}\, \left( \text{LLM} (\text{RoPE}, \mathbf{X})\right), \ \text{where} \;
\\
\scalebox{0.93}{
$\underset {\begin{subarray}{l} i=0,\cdot\cdot,\frac{d}{2}-1;\\ n\in[ 0,|\mathbf{x}|);\end{subarray}}{\text{RoPE}_(n)}= {\left[\cdots, \cos\left(\mathbb{I}(\hat{\lambda}_{i}, \hat{n})\frac{n}{\beta^i}\right), \sin\left(\mathbb{I}(\hat{\lambda}_{i}, \hat{n})\frac{n}{\beta^i}\right), \cdots\right] }$}
\\
	\text{where} \; \mathbb{I}(\hat{\lambda}_{i},\hat{n}) = \begin{cases}
	1&\mbox{\; if $n<\hat{n}$}\\
	\frac{1}{\lambda_{i}}&\mbox{\; if $n\geq\hat{n}$}
\end{cases}
	\end{gathered}
\end{equation}

In this approach, a collection of rescaling coefficients, $\mathbb{I}(\hat{\lambda}_{i},\hat{n})$, is introduced to address both types of non-uniformity present in the system. Here, $\hat{\lambda_i}$ refers to the variation along the RoPE dimensions, while $\hat{n}$ indicates discrepancies related to token positions. The function $\mathbb{I}(\hat{\lambda}_{i},\hat{n})$ specifically adjusts the rotational angle in the $i^{th}$ RoPE dimension, utilizing $\hat\lambda_{i}$ as the scaling parameter and $\hat{n}$ as the position threshold. For the initial $\hat{n}$-1 tokens, $\hat\lambda_{i}$ is inactive, so the conventional RoPE rotation, $\frac{n}{\beta^i}$, remains in use. Once token positions reach $n\ge\hat{n}$, the rescale factor $\hat\lambda_{i}$ is activated and applied.

With a specified context window size $L'$, the task involves determining the set of optimal rescale coefficients ($\mathbb{I}(\hat{\lambda}_{0},\hat{n})$, $\mathbb{I}(\hat{\lambda}_{1},\hat{n})$, ... $\mathbb{I}(\hat{\lambda}_{i},\hat{n})$, ...) corresponding to each RoPE dimension from the $1^{st}$ up to the $d^{th}$. By applying these adjustments to the RoPE of the target $\text{LLM}$, the aim is to minimize the next token prediction loss, denoted by $\mathcal{L}$ (which reflects perplexity), when processing input sequences $\mathbf{X}$ consisting of $L'$ tokens.

\subsection{Searching the Non-uniform Position Interpolation}

To address the challenge outlined in Eq.~\ref{eq:problem}, we present a straightforward and powerful approach that identifies the best RoPE rescale factors, thereby leveraging the diverse multidimensional patterns found in position embedding.

\textbf{Search space}. To capture the two non-uniform characteristics, we establish an extensive search space. The structure of this space is detailed in Table~\ref{tbl:searchspace}. In our formulation, we permit the optimization of an individualized rescale factor for every dimension within the RoPE embedding. For simplicity in the search space construction, we focus on searching over $\lambda_i$ and $\hat{n}$, rather than directly optimizing $\mathbb{I}(\hat{\lambda}_{i},\hat{n})$, given that $\hat{\lambda}_{i}=1/\lambda_i$. Table~\ref{tbl:searchspace} demonstrates that $\lambda_i$ ranges from a minimum of 1.0 (corresponding to straightforward extrapolation) up to a maximum of $s\times1.25$ (which enables interpolation beyond PI), with increments of 0.01. Here, $s$ signifies the desired ratio of the context window's extension.

The parameter $\hat{n}$ determines how many leading token positions preserve the original RoPE embedding, foregoing position interpolation. In practice, we permit $\hat{n}$ to be selected from the set \{0, 1, 2, 4, 8, 12, 16, 20, 24, 28, 32, 64, 128, 256\}. Setting $\hat{n}=0$ implies that every token position employs the rescale factors obtained from the search process.

\textbf{Evolution-based search}. The search space outlined in Table~\ref{tbl:searchspace} includes a vast array of positional interpolation configurations, presenting considerable challenges for effective search strategies. For instance, with an extension rate of $s=4\times$, there exist $400^{128/2}\times14 = 4\times10^{167}$ possible selections. As the extension ratio increases, the number of possible combinations grows exponentially. To mitigate this complexity, we employ evolution search~\cite{guo2020single}, supplemented by two optimization strategies designed to significantly enhance search performance. The complete search process is depicted in Algorithm~\ref{alg:search}.

\textit{Optimized initial population generation}. In place of purely random initialization for the population of $P$ rescale factors, the initial population explicitly includes the three RoPE rescale factors associated with PI, NTK, and YaRN as individual members. The other $P-3$ individuals are generated by applying random mutations, with probability $p$, to these three rescale factors.

\textit{Monotonically non-decreasing constraint}. Once the initial population is created, we evaluate LLM perplexity for every individual by applying the specific RoPE rescale factors to the target LLM and then measuring perplexity on input $\mathbf{X}$. The top-k candidates based on this score are selected as parents for the subsequent evolutionary steps. Nonetheless, due to the extensive search space, straightforward mutation and crossover operations may frequently yield suboptimal candidates, resulting in many redundant perplexity evaluations. This inefficiency becomes especially pronounced as $L'$ increases, since each perplexity assessment involves computationally expensive inference.

To mitigate this issue, we enforce a non-decreasing monotonic constraint on the sampled RoPE rescaling coefficients: $\lambda_i\le \lambda_{i+1}$. We only utilize RoPE configurations that comply with this rule for perplexity analysis within LLMs, thereby considerably decreasing the computational burden of the search process. Concretely, we stipulate that $\lambda_i$ must grow monotonically as a function of the RoPE dimension (where $i=0,\ldots,63$). This requirement for monotonicity across dimensions draws its justification from NTK theory~\cite{jacot2018neural,tancik2020fourier,ntk}, which posits that dimensions associated with higher frequency (lower indices) need minimal interpolation (that is, smaller $\lambda_i$ values), while dimensions corresponding to lower frequencies (higher indices) permit greater interpolation (thus, larger $\lambda_i$ values).

\begin{algorithm}[t]
	\small
	\caption{The search algorithm}
	\textbf{Input:} target LLM, input samples $\mathbf{X}$, population size $P$, mutation size $N_1$, crossover size $N_2$, maximum iterations $\mathcal{T}$, mutation probability $p$\\

	\begin{algorithmic}[1]
		\label{alg:search}
		\STATE Top-k $\,\gets\,$ $\phi$;
		\STATE $\text{P}_0 \gets$ \textit{Initialize\_population\_with\_optimization} ($P$, $\mathbf{X}$, $p$);
		\FOR{$i = 1$ \textbf{to} $\mathcal{T}$}
		\STATE \textit{Compute\_perplexity} (LLM, $\text{P}_{i-1}$, $\mathbf{X}$);
		\STATE Top-k $\gets$ \textit{Update\_Topk} (Top-k, $\text{P}_{i-1}$);
		\STATE $\text{P}_{mutation} \gets$ \textit{Mutation\_with\_mono\_constraint} (Top-k, $N_1$, $p$);
		\STATE $\text{P}_{crossover} \gets$ \textit{Crossover\_with\_mono\_constraint} (Top-k, $N_2$);
		\STATE $\text{P}_i \gets \text{P}_{mutation} \cup \text{P}_{crossover} \cup$ Top-k;
		\ENDFOR
		\STATE Output the individual within Top-k that achieves the lowest perplexity;
	\end{algorithmic}
\end{algorithm}

\textbf{8$\times$ expansion without additional training}. Through evolutionary search, our approach efficiently uncovers non-uniform RoPE rescale factors, maintaining crucial axes and positional encoding to reduce information degradation from interpolation. As shown in Fig.\ref{fig:non-ft}, our technique successfully broadens the context window of LLaMA2 from 4k to 32k without requiring any fine-tuning. Conversely, alternative strategies like PI and the non-uniform versions of NTK and YaRN result in a rapid increase in perplexity when extending the window beyond 2$\times$.

\subsection{Extending LLM Context Window to 2048K}
\label{sec:extendto2048k}

\textbf{Progressive extension to 2048k}. In this section, we propose an approach for increasing the context window of pre-trained LLMs from the standard 4k up to beyond 2048k. As evidenced previously, our non-uniform positional interpolation can provide an 8$\times$ context window increase without the need for fine-tuning. When aiming for much larger expansions (for example, 512$\times$), however, fine-tuning becomes indispensable. One possible strategy involves selecting appropriate RoPE rescaling parameters for the target 2048k length prior to initiating fine-tuning. Nonetheless, this process is hindered by the extremely high computational costs it entails. Furthermore, our observations indicate that performing effective fine-tuning on LLMs with such a significant extension ratio is quite difficult (refer to Appendix).

Fortunately, LongRoPE demonstrates robust performance with both pretrained and further fine-tuned extended LLMs. Accordingly, we propose an efficient and incremental approach that reaches the intended 2048k context length using only 1k fine-tuning steps, all within a maximum training sequence length of 256k.

$\Diamond$ \textit{Exploring expansion of pre-trained LLMs to 256k using LongRoPE search}. Using LLaMA2 as a reference model, we explore configurations targeting context windows of 128k and 256k. At these sizes, the model’s context length is increased by factors of 32$\times$ and 64$\times$, correspondingly.

$\Diamond$ \textit{Fine-tuning to 256k}. In this stage, we adapt the pre-trained LLM to support a context window of 256k tokens. Our approach involves initially fine-tuning LLaMA2 with RoPE rescaled factors set to 128k for 400 steps. After completing this phase, we update the rescaled factors to 256k in the resulting checkpoint and proceed with 600 more fine-tuning steps. This two-stage process offers greater efficiency compared to fine-tuning for 256k context length in a single step.

$\Diamond$ \textit{LongRoPE search enables scaling fine-tuned extended LLMs to 2048k}. In the last step, we conduct an additional search on the LLM that has already been fine-tuned for 256k-length sequences. This process increases the context window to a massive 2048k, all accomplished without the need for further fine-tuning. This achieves a total extension factor of $512\times$.

\textbf{Restoration within restricted context windows}.  When expanding the context window to a remarkable 2048k length, a decline in performance is observed within the initial context window. This drawback is recognized in the literature concerning positional interpolation~\cite{pi}, which compresses the position embeddings in the original context into a reduced space in higher dimensions, thus impairing the model's language comprehension capabilities. Specifically, at a 512$\times$ expansion rate, the positions in the original 4k context window are forced into an unusually dense configuration.

To address this issue, we conduct an additional evolutionary search on the extended LLM, specifically tuning the RoPE rescale factors for shorter context lengths (such as 4k and 8k). For these shorter sequences, we restrict the upper bound of the searched $\lambda$, as they necessitate less positional interpolation. At inference time, the LLM automatically applies the appropriate RoPE rescale factors based on the input context length.

\section{LongRoPE}

\label{sec:method}
Drawing inspiration from the aforementioned results, we introduce LongRoPE. This method incorporates an optimized search procedure specifically designed to harness both identified non-uniformities. Utilizing this approach, LongRoPE is able to increase the LLM context window capacity to exceed 2 million tokens.

\subsection{Problem Formulation}

These two forms of non-uniformity expand the solution space significantly, thereby complicating the optimization process. To tackle this issue, we conceptualize the optimization challenge of multidimensional non-uniform position interpolation as a search problem.

For a LLM targeting a  context window size of $L'$ and lengthy input documents $\mathbf{X}$, where each $\mathbf{x}\in\mathbf{X}$ surpasses $L'$ in token length, we denote the original rotary angle of the $i^{th}$ dimension in RoPE embedding at token position $n$ as  $\frac{n}{\beta^i}$. The optimization problem is then formulated as follows:

\begin{equation}
\begin{gathered}
		\label{eq:problem}

\underset {\mathbf{x}\in\mathbf{X}; \, |\mathbf{x}|\ge L'}{ \operatorname {arg\,min} } \,  \mathcal{L}\, \left( \text{LLM} (\text{RoPE}, \mathbf{X})\right), \text{where} 
\\
\scalebox{0.93}{
$\underset {\begin{subarray}{l} i=0,\cdot\cdot,\frac{d}{2}-1;\\ n\in[ 0,|\mathbf{x}|);\end{subarray}}{\text{RoPE}_{(n)}}= {\left[\cdots,\cos\left(\mathbb{I}(\hat{\lambda}_{i}, \hat{n})\frac{n}{\beta^i}\right),  \sin\left(\mathbb{I}(\hat{\lambda}_{i}, \hat{n})\frac{n}{\beta^i}\right),\cdots\right] }$}\\
	\text{where} \quad \mathbb{I}(\hat{\lambda}_{i},\hat{n})=\begin{cases}
	1 &\text{\; if $n < \hat{n}$}\\
	{\frac{1}{\lambda}_{i}} &\text{\; if $n \geq \hat{n}$}
\end{cases}
	\end{gathered}
\end{equation}

We define a group of rescale factors $\mathbb{I}(\hat{\lambda}_{i},\hat{n})$ to address two types of non-uniformity in the system. Here, $\hat{\lambda}_{i}$ corresponds to irregularities across RoPE dimensions, while $\hat{n}$ captures non-uniformity at certain token positions. The function $\mathbb{I}(\hat{\lambda}_{i},\hat{n})$ adjusts the rotation angle for the $i^{th}$ dimension in RoPE by applying the rescale factor $\hat{\lambda}_{i}$ when token positions meet or exceed the threshold $\hat{n}$. For the first $\hat{n}-1$ positions, rescaling does not occur and the standard RoPE angle $\frac{n}{\beta^i}$ is retained. Once $n$ reaches $\hat{n}$ or beyond, the adjustment via the rescale factor takes place.

Assuming a desired context window of size $L'$, the task is to determine the most suitable rescale factors ($\mathbb{I}(\hat{\lambda}_{0},\hat{n})$, $\mathbb{I}(\hat{\lambda}_{1},\hat{n})$, ..., $\mathbb{I}(\hat{\lambda}_{i},\hat{n})$...) corresponding to each RoPE dimension from the $1^{st}$ up to the $d^{th}$. With these adjusted $\text{RoPE}$ scalings, the target $\text{LLM}$ is expected to yield the lowest possible next token prediction loss, $\mathcal{L}$ (reflecting model perplexity), over input sequences $\mathbf{X}$ whose lengths are $L'$.

\subsection{Searching the Non-uniform Position Interpolation}

To address the challenge posed in Eq.~\ref{eq:problem}, we propose a straightforward but powerful approach that identifies the ideal RoPE rescale factors, thereby leveraging the complex, non-uniform nature of multidimensional position embeddings to their fullest potential.

\textbf{Search space}. We establish a broad search space to capture the two forms of non-uniformity. The scope of this search is summarized in Table~\ref{tbl:searchspace}. More precisely, we permit the optimization of individual rescale factors for each dimension within the RoPE embedding. In order to streamline the search space, the search procedure focuses on tuning $\lambda_i$ and $\hat{n}$, rather than directly optimizing $\mathbb{I}(\hat{\lambda}_{i},\hat{n})$, with the relation $\hat{\lambda}_{i}=1/\lambda_i$. According to Table~\ref{tbl:searchspace}, the range for $\lambda_i$ spans from a lower bound of 1.0 (representing plain extrapolation) to an upper bound of $s\times1.25$ (corresponding to enhanced interpolation compared to PI), incremented by 0.01, where $s$ denotes the target ratio for extending the context window.

The parameter $\hat{n}$ specifies how many of the initial token positions are kept with their original RoPE embeddings, foregoing any form of position interpolation. In practical terms, we explore values for $\hat{n}$ from the set \{0, 1, 2, 4, 8, 12, 16, 20, 24, 28, 32, 64, 128, 256\}. If $\hat{n}=0$, this means that the rescale factors determined through search are applied to every token position.

\textbf{Evolution-based search}. As shown in Table~\ref{tbl:searchspace}, the scope of our search encompasses a vast array of positional interpolation configurations, making systematic navigation through the space highly challenging. For instance, selecting $s=4\times$ for extension yields $400^{128/2}\times14 = 4\times10^{167}$ possible arrangements. This exponential increase in the number of candidates is amplified with higher extension ratios. To tackle this complexity, we implement evolution search~\cite{guo2020single} alongside two key optimization strategies designed to markedly improve search efficiency. The entire search workflow is presented in Algorithm~\ref{alg:search}.

\textit{Optimized initial population generation}. Rather than relying solely on random initialization for the population of $P$ rescale factors, the approach includes the three RoPE rescale factors associated with PI, NTK, and YaRN directly as members of the initial population. The other $P-3$ individuals are created by applying random mutations to these three rescale factors, each with a mutation probability of $p$.

\textit{Monotonically non-decreasing constraint}. Once the initial population is formed, we evaluate the LLM perplexity of each candidate by applying their specific RoPE rescale factors to the target LLM and determining the perplexity of the input $\mathbf{X}$. The best-performing $k$ individuals are then chosen as parents for the next evolutionary stage. Nevertheless, due to the enormous search space, unsophisticated mutation and crossover processes may generate suboptimal candidates, resulting in excessive and redundant perplexity evaluations. This inefficiency is further magnified when $L'$ is large, as each perplexity computation incurs significant computational overhead.

To mitigate this issue, we enforce a monotonicity constraint such that the sampled RoPE rescaling factors are non-decreasing: $\lambda_i\le \lambda_{i+1}$. Only those RoPE configurations that adhere to this requirement are utilized for perplexity assessment in LLMs, thereby streamlining the search process considerably. Concretely, we stipulate that $\lambda_i$ must grow monotonically with respect to the RoPE dimension index ($i = 0, ..., 63$). This pattern of monotonicity across dimensions is motivated by NTK theory~\cite{jacot2018neural,tancik2020fourier,ntk}, which posits that dimensions corresponding to higher frequencies (i.e., lower indices) necessitate less interpolation and thus smaller $\lambda_i$, while dimensions associated with lower frequencies (higher indices) are amenable to greater interpolation, justifying larger $\lambda_i$ values.

\begin{algorithm}[t]
	\small
	\caption{The search algorithm}
	\textbf{Input:} target LLM, input samples $\mathbf{X}$, population size $P$, mutation size $N_1$, crossover size $N_2$, maximum iterations $\mathcal{T}$, mutation probability $p$\\

	\begin{algorithmic}[1]
		\label{alg:search}
		\STATE Initialize Top-k as an empty set ($\phi$);	
		\STATE Create initial population $\text{P}_0$ via \textit{Initialize\_population\_with\_optimization} ($P$, $\mathbf{X}$, $p$); 
		\FOR{$i$ from $1$ to $\mathcal{T}$}
		\STATE Evaluate perplexity for $\text{P}_{i-1}$ using the LLM and $\mathbf{X}$ through \textit{Compute\_perplexity};
		\STATE Update Top-k selection using \textit{Update\_Topk} with current Top-k and $\text{P}_{i-1}$;
		\STATE Generate mutated population $\text{P}_{mutation}$ via \textit{Mutation\_with\_mono\_constraint} on Top-k, with parameters $N_1$ and $p$;
		\STATE Produce crossover population $\text{P}_{crossover}$ via \textit{Crossover\_with\_mono\_constraint} applied to Top-k and $N_2$;
		\STATE Assemble next population $\text{P}_i$ as the union of $\text{P}_{mutation}$, $\text{P}_{crossover}$, and Top-k;
		\ENDFOR
		\STATE Output the member in Top-k with the minimal perplexity;
	\end{algorithmic}
\end{algorithm}

\textbf{8$\times$ extension without fine-tuning}. Leveraging evolutionary search, our approach discovers optimal non-uniform RoPE rescale factors, strategically maintaining crucial dimensions and positional encoding to reduce information loss resulting from interpolation. As illustrated in Fig.\ref{fig:non-ft}, this technique enables LLaMA2 to achieve a context window expansion from 4k to 32k tokens without any fine-tuning. In comparison, current approaches—including PI as well as non-uniform NTK and YaRN—show marked increases in perplexity when the context window is doubled.

\subsection{Extending LLM Context Window to 2048K}
\label{sec:extendto2048k}

\textbf{Progressive extension to 2048k}. In this section, we present our approach for increasing the context window of pre-trained LLMs from the standard 4k tokens up to beyond 2048k tokens. As shown previously, our non-uniform positional interpolation technique enables an 8$\times$ enlargement of the context window without requiring any fine-tuning. For more substantial expansions—such as reaching a 512$\times$ extension—fine-tuning becomes essential. A common strategy involves identifying suitable RoPE rescaling factors corresponding to the desired 2048k context size, followed by a fine-tuning phase. However, this approach encounters significant difficulties due to the extremely high computational costs involved in training. Additionally, our observations indicate that fine-tuning LLMs effectively at such extensive extension ratios remains a considerable challenge (refer to Appendix).

Fortunately, LongRoPE demonstrates efficacy with both the pre-trained and the fine-tuned versions of extended LLMs. Consequently, we propose a resource-efficient, incremental approach capable of reaching the target 2048k by utilizing only 1k fine-tuning iterations within a maximum training sequence length of 256k.

$\Diamond$ \textit{Utilizing LongRoPE search to scale pre-trained LLMs up to 256k}. For instance, using LLaMA2, we apply LongRoPE search targeting context lengths of 128k and 256k tokens. At this phase, the expansion factors are 32$\times$ and 64$\times$, respectively.

$\Diamond$ \textit{Fine-tuning to 256k}. In this phase, the pre-trained LLM is adjusted to support a context window of 256k. The process begins by fine-tuning LLaMA2 for 400 steps utilizing RoPE rescaled factors calibrated for 128k. Subsequently, these RoPE rescaled factors are swapped for those corresponding to 256k on the already trained checkpoint, followed by an extra 600 steps of fine-tuning. This approach demonstrates greater efficiency compared to immediately fine-tuning at 256k.

$\Diamond$ \textit{Scaling a previously fine-tuned extended LLM to 2048k with a LongRoPE search.} In the last step, we apply an additional search process to the LLM already fine-tuned at 256k-length. This approach enables expansion to an extensive context window size of 2048k, achieved without additional fine-tuning. The overall increase in window length corresponds to a $512\times$ extension.

\textbf{Restoration of Performance in Short Contexts}.  Expanding the context window to a substantial length of 2048k leads to diminished performance within the initial, smaller context range. This challenge is characteristic of positional interpolation methods~\cite{pi}, which confine position embeddings corresponding to the original window into a much smaller segment of the overall positional space. This spatial compression adversely influences the model’s abilities. Specifically, when applying a 512$\times$ expansion factor, the original 4k context window positions are densely packed together, exacerbating the degradation.

To address this issue, we conduct an additional evolution search on the extended LLM, specifically targeting the optimization of RoPE rescale factors at shorter context lengths (such as 4k and 8k). For these reduced sequence lengths, the search space for the maximal $\lambda$ is constrained, as the demand for positional interpolation is diminished. At inference time, the LLM updates the associated RoPE rescale factors adaptively.

 \section{Experiments}

\subsection{Setup}
\textbf{Evaluation Tasks and models}. We utilize LongRoPE with both LLaMA2-7B and Mistral-7B, assessing outcomes based on three criteria: (1) the perplexity scores of large language models with expanded contexts when processing lengthy texts; (2) performance on the Passkey retrieval task, which tests the model's capability to locate a specific passkey amidst extensive irrelevant information; and (3) results on conventional LLM benchmarks restricted to a 4096 token context window.

\textbf{Fine-tuning}. For LLaMA2, we adopt a learning rate of $2 \times 10^{-5}$, applying linear decay, with a total batch size of 32. The model undergoes 400 fine-tuning steps on the Redpajama~\cite{redpajama} dataset, which is divided into segments of 128k tokens and includes both BOS and EOS markers at the boundaries. Following this initial phase, we conduct an additional 600 training steps using the resulting checkpoint, extending the context window to 256k tokens. Training with 128k context is carried out on a cluster of 8 A100 GPUs utilizing a distributed training infrastructure~\cite{cube}, while scaling up to 256k context necessitates 16 A100 GPUs. For Mistral, the learning rate remains fixed at $1 \times 10^{-6}$ with a global batch size of 64. Both 128k and 256k context models replicate the YaRN~\cite{yarn} configuration, completing 400 training steps on the Long-Data Collections from Together Computer~\cite{mistraldata}, using sequences of 16k tokens. This process employs 4 A100 GPUs.

\textbf{Search}. When the target window size does not exceed 256k, we set the parameters as follows: $P$=64, $N_1$=$N_2$=16, $p$=0.3, $\mathcal{T}$=40, and in each generation, the top-32 candidates are chosen for mutation and crossover operations. Perplexity is computed using 5 randomly chosen samples from the PG19 validation set, ensuring each meets or exceeds the target context length. For window sizes greater than 512k, we reduce the population as well as the numbers for mutation and crossover by half. In this case, perplexity evaluation uses 3 random samples from the Pile-Books3~\cite{pile} validation set.

\textbf{Baselines}. To achieve a 2048k context window, we fine-tuned models with initial context sizes of 128k and 256k. As a result, we obtain LongRoPE-2048k (ft=128k) for LLaMA2 and LongRoPE-2048k (ft=256k) for Mistral. We evaluate these four models against leading methods for extending context windows, including several open-source LLMs that have undergone fine-tuning following positional interpolation via PI, NTK, and YaRN. These reference baselines comprise Together-32k~\cite{together}, Code LLaMA~\cite{codellama}, LongLoRA-full-FT-100k~\cite{longlora}, YaRN-LLaMA, and YaRN-Mistral~\cite{yarn}.

\subsection{Main Results}

\label{sec:mainresult}
\textbf{Language modeling on extended sequences up to 256k.} Our analysis commences with a comparison to state-of-the-art long-context LLMs over evaluation sequences reaching 256k tokens. To establish generalizability, we utilize two datasets: the test splits from Proof-pile~\cite{pg19} and PG19~\cite{pile}. Perplexity scores are obtained across multiple context lengths, employing a 256-token sliding window. For PG19, the entire test split containing 100 documents is evaluated, while for Proof-pile, in accordance with YaRN~\cite{yarn}, we randomly draw 10 samples, each measuring at least 128k tokens in length.

Tables~\ref{tbl:proof-pile} and~\ref{tbl:pg19} present a comparison of perplexity scores for LLaMA2 and Mistral models, which have been extended using various interpolation techniques, across the Proof-pile and PG19 datasets. Two main findings emerge from this analysis: \textbf{(1)} the perplexity of our extended models consistently declines as the evaluation length increases from 4k to 256k, demonstrating an improved ability to utilize longer context effectively; \textbf{(2)} even when handling a context window that is 16$\times$ larger—a scenario in which it is typically difficult to preserve performance for shorter sequences—our LongRoPE-2048k variants achieve better results than current leading baselines within the 256k context length range.

\textbf{Modeling long sequences beyond 2000k tokens}. To assess performance on very lengthy texts, we employ the Books3~\cite{pile} dataset. To streamline evaluation, we randomly sample 20 books, each with a length greater than 2048k, applying a sliding window approach of size 256k.

Table~\ref{tbl:books3} demonstrates that LongRoPE is able to expand the context window of LLaMA2-7B and Mistral-7B up to 2048k, while maintaining perplexity scores that match or surpass baseline performance at shorter context lengths (8k-128k). In comparing results between LLaMA2 and Mistral at the 2048k length, substantial performance discrepancies emerge. Mistral exhibits superior results against baselines for shorter contexts; however, its perplexity rises above 7 once the window surpasses 256k tokens. LLaMA2 generally displays anticipated behavior: perplexity reliably improves (decreases) as the window grows, although slight increases are noted at 1024k and 2048k. Notably, for LLaMA2, LongRoPE-2048k yields better outcomes when fine-tuned with a 256k sequence length rather than 128k, attributable to a more modest secondary extension ratio (8$\times$ compared to 16$\times$). Conversely, Mistral achieves optimal performance at a fine-tuning length of 128k. This difference is primarily because, following the YaRN protocol, both the 128k and 256k Mistral fine-tuning experiments employ a 16k training length, which limits Mistral’s capacity to further expand its context window post fine-tuning.

\textbf{Passkey retrieval}. Next, we examine the usable context window length for generation tasks. To do so, we utilize a synthetic passkey retrieval benchmark, as introduced by~\cite{passkey}. In this setup, the model must locate and extract a randomly assigned passkey (a sequence of five digits) embedded within an extended document. The full prompt structure is provided in the appendix. We conduct ten runs of the passkey retrieval experiment, with the passkey inserted at a randomly chosen position that is uniformly sampled from the entire context window used for evaluation.

Fig.~\ref{fig:passkey} presents retrieval accuracy results alongside baseline models. The accuracy of prior models decreases sharply to zero when the sequence length exceeds 128k. By contrast, LongRoPE-LLaMA2-2048k (ft=256k) demonstrates impressive performance, sustaining high retrieval accuracy ($\ge$90\%) across the entire range from 4k to 2048k tokens—even when tasked with extracting a passkey from sequences at the million-token scale. Similarly, LongRoPE-Mistral-2048k (ft=128k) achieves perfect accuracy up to 1800k tokens, only dropping to 60\% at 2048k. This outcome is consistent with observations in Table~\ref{tbl:books3}, where perplexity exhibits a moderate increase at 2048k.

\textbf{Evaluation on standard benchmarks within the initial context window}. LongRoPE-2048k models are assessed in their native context window using the Hugging Face Open LLM Leaderboard~\cite{huggingfaceboard} under both zero-shot and few-shot paradigms. Specifically, we employ 25-shot ARC-Challenge~\cite{allenai:arc}, 10-shot HellaSwag~\cite{zellers2019hellaswag}, 5-shot MMLU~\cite{mmlu}, and 0-shot TruthfulQA~\cite{lin2021truthfulqa}.

According to Table~\ref{tbl:commonsense}, the models presented demonstrate similar performance to the original benchmark, which was tailored for a shorter context window. They also surpass the original Mistral model on TruthfulQA by an improvement of +0.5\%. While LongRoPE-LLaMA2-2048k, fine-tuned at 256k, experiences somewhat greater reduction in effectiveness, its performance across the majority of tasks continues to fall within acceptable margins.

\subsection{Ablation Results}

\textbf{Assessing the Effectiveness of the Second Positional Interpolation}. Within our progressive extension framework, we implement a search algorithm to apply an additional non-uniform positional interpolation to the fine-tuned, extended LLMs. To examine its impact, we perform evaluations on the fine-tuned LLaMA2-256k model. This model is further extended to 512k, 1024k, and 2048k positions using both PI and YaRN techniques. As observed in Table~\ref{tbl:secondsearch}, the use of our non-uniform positional interpolation enables the model to maintain stable perplexity scores. In comparison, both PI and YaRN demonstrate rapidly increasing perplexity as the extension ratio grows.

\textbf{Recovery performance at reduced context lengths.} To address declines in performance at lower context lengths, we recalibrate the RoPE factors for LongRoPE-2048k using our search strategy. In particular, we limit the maximum scale factors allowed during the search process to reduce the degree of interpolation applied at shorter context lengths such as 4k and 8k. 
 Table~\ref{tbl:shortperformance} presents a comparison of perplexity for LongRoPE-LLaMA2-2048k on the Proof-pile dataset at 4k and 8k context windows, as well as the mean LLM benchmark accuracy. The findings indicate notable gains in performance at these shorter context lengths.

\textbf{Analysis on the two forms of non-uniformities}. In this section, we investigate the influence of the two non-uniformities through an ablation study to assess their respective contributions to model performance. Two separate experimental setups are employed: (i) expanding LLaMA2-7B to shorter sequence lengths of 16k and 32k, comparing the approaches of PI, tuning only the RoPE dimension, and tuning both non-uniformities; (ii) increasing the context length of our fine-tuned LLaMA2-256k model to 2048k, using analogous strategies. Perplexity metrics are computed without additional fine-tuning. According to Table~\ref{tbl:searchdimension}, adjusting the RoPE dimension non-uniformity yields a substantial perplexity reduction compared to PI's method of linear interpolation. Introducing non-uniformity in token positions notably benefits results at 16k and 32k, but its effect diminishes at 2048k—potentially due to the excessive sequence length, where relying solely on the initial token distribution without interpolation ceases to be effective. The investigation of this phenomenon is deferred to future research directions.

\section{Related Works}

Besides position interpolation techniques, this section examines alternative strategies found in related literature.

\textbf{Retrieval-based} methods incorporate external memory mechanisms to store extensive contextual information and leverage retrieval modules to access pertinent documents during inference~\cite{longllama,wang2023augmenting,borgeaud2022improving}. These approaches usually require significant architectural alterations to large language models. In contrast, our method remains comparatively lightweight by restricting modifications to minor changes in positional embeddings. Furthermore, our approach extends to a wider range of long-context applications beyond mere retrieval, including long document summarization and few-shot learning.

\textbf{Attention-based context window extensions}. In addition to approaches that interpolate positional embeddings, other works have expanded input context capacity within the constraints of the native LLM context window by altering the attention mechanism~\cite{han2023lminfinite, streamingllm,parallelwindow}. Central to these methods is the introduction of new attention masks aimed at addressing the issue of increased computational overhead from novel positions. These techniques serve as a complement to positional interpolation strategies.

\textbf{Fine-tuning based approaches} investigate strategies to efficiently adapt pre-trained LLMs to longer sequences by adjusting their position embeddings. Methods such as Code LLaMA~\cite{codellama}, LLaMA2 Long~\cite{xiong2023effective}, and ScaledRoPE~\cite{liu2023scaling} adopt extremely large base values for RoPE during fine-tuning at the intended context length. The approach presented in this work enables compatibility with multiple target sequence lengths and extends beyond 2M tokens. Recently, due to the heavy GPU requirements for fine-tuning at very large contexts (exceeding 128k tokens), methods like LongLoRA~\cite{longlora} and PoSE~\cite{zhu2023pose} have been introduced to address resource consumption. The technique we propose remains complementary and independent of these optimization strategies for efficient fine-tuning.

\section{Conclusion}

This paper introduces LongRoPE, a novel approach that dramatically enhances the context length of LLMs to 2048k tokens—far surpassing previous limits—while preserving their performance with standard, shorter context windows. The method utilizes two distinct forms of non-uniformity in RoPE positional embeddings, identified via a highly efficient evolutionary search. This provides two major advantages: effective initialization for subsequent fine-tuning and an 8$\times$ increase in context window length even without fine-tuning. Leveraging these benefits, we put forward a gradual extension strategy, using LLMs fine-tuned on 256k-length contexts, which allows us to achieve a context window of 2048k without any additional fine-tuning steps. Comprehensive empirical analysis demonstrates the success of LongRoPE. We anticipate that LongRoPE-2048k will facilitate a wide range of long context applications and motivate new directions for future studies.

\section*{Broader Impacts} The aim of this paper is to contribute to the progress of Machine Learning research. Our findings may have a range of effects on society, but we do not identify any particular outcomes that warrant emphasis at this time.

	\balance

\end{document}